\topic{附录 B\quad 60 套试卷逐卷题型对照}
\subsection*{B.1 常规初赛：主要统计样本}
{\small\renewcommand{\arraystretch}{1.18}
\begin{longtable}{P{2.8cm}P{1.8cm}P{10.8cm}}
\toprule 卷别 & 真题 PDF 页 & 各题主标签（依原题序）\\\midrule\endhead
1 初赛 & 5--7 & 题序：G、L3、L1、A4、A4、A3、A2\\
2 初赛 & 12--14 & 题序：A1、L3、A2、A3、G、L4、A3、A5\\
3 初赛 & 18--21 & 题序：G、A3、L3、A2、A1、L1、A5\\
4 初赛 & 25--27 & 题序：G、G、A5、L4、A4、A2、L4\\
5 初赛 & 31--33 & 题序：G、L5、A1、A2、A3、L3\\
6 初赛 & 42--44 & 题序：G、A3、L3、A2、L3、A1\\
7 初赛 & 53--55 & 题序：G、L1、L3、A1、A4、A3\\
8 初赛 & 64--66 & 题序：G、L1、L2、A1、A2、A3\\
9 初赛 & 76--78 & 题序：G、A4、L1、L2、A1、A3\\
10 初赛 & 87--89 & 题序：G、L4、L5、A2、A4、A2\\
11 A 初赛 & 100--102 & 题序：G、A2、L1、L4、A1、A2\\
11 B 初赛 & 103--105 & 题序：G、A5、A1、L3、L1、A2、A4\\
12 A 初赛 & 116--118 & 题序：G、A1、L4、L5、A1、A3\\
12 B 初赛 & 119--120 & 题序：G、A3、L5、L3、A5、A3\\
13 A 初赛 & 132--134 & 题序：G、A2、L5、L4、A5、A2\\
13 B 初赛 & 135--136 & 题序：G、A2、L5、L2、A4、A2\\
14 A 初赛 & 152--154 & 题序：G、A1、L3、L5、A3、A6\\
14 B 初赛 & 155--156 & 题序：G、A1、L3、L1、A4、A2\\
15 A 初赛 & 177--179 & 题序：G、A2、L3、L2、A6、A1\\
15 B 初赛 & 180--182 & 题序：G、A2、L3、L2、A2、A5\\
16 A 初赛 & 193--195 & 题序：G、A5、L2、L3、A4、A1\\
16 B 初赛 & 196--198 & 题序：G、A5、L2、L2、A1、A1\\
17 A 初赛 & 208--210 & 题序：G、A3、L2、L3、A2、A3\\
17 B 初赛 & 211--213 & 题序：G、A3、L2、L2、A1、A2\\
\bottomrule\end{longtable}}

\clearpage
\subsection*{B.2 补赛：另列作为稳健性检查}
{\small\renewcommand{\arraystretch}{1.18}
\begin{longtable}{P{2.8cm}P{1.8cm}P{10.8cm}}
\toprule 卷别 & 真题 PDF 页 & 各题主标签（依原题序）\\\midrule\endhead
13 A 补赛 & 137--139 & 题序：G、A5、L3、L1、A1、A4\\
13 B 补赛 & 140--141 & 题序：G、A5、L3、L1、A3、A4\\
14 A 补一 & 157--159 & 题序：G、A3、L1、L1、A1、A5\\
14 B 补一 & 160--161 & 题序：G、A3、L1、L2、A2、A2\\
14 A 补二 & 162--164 & 题序：G、A6、L5、L2、A2、A2\\
14 B 补二 & 165--166 & 题序：G、A6、L5、L1、A1、A4\\
\bottomrule\end{longtable}}

\para{补赛提供的额外提醒}
第 14 届补二的集合与可导构造说明，初赛仍可能出现较抽象的题；但把这些集中在同届补赛的题位单列后，不宜据此大幅挤压长期稳定的数列、导数、积分、线性结构和几何训练。第 14 届补一的含参积分，则与第 16 届常规初赛的方法有清楚联系，适合作为掌握方法后的变式。
\subsection*{B.3 早期未分组决赛}
{\small\renewcommand{\arraystretch}{1.18}
\begin{longtable}{P{2.8cm}P{1.8cm}P{10.8cm}}
\toprule 卷别 & 真题 PDF 页 & 各题主标签（依原题序）\\\midrule\endhead
1 决赛 & 8--11 & 填：A5、A2、A1、G\newline 答：A1、A4、A2、A3、G、L4、L3\\
2 决赛 & 15--17 & 题序：G、A3、A4、L1、A6、L3\\
3 决赛 & 22--24 & 题序：G、A3、A2、A1、L1、L3、A5\\
4 决赛 & 28--30 & 题序：G、A1、L2、A2、L4、A3\\
\bottomrule\end{longtable}}

\para{为什么不混入初赛频率}
早期决赛的计算、证明和综合题形式与后期分组不同。本表保留它们供回查；分析初赛优先级时仍用常规初赛为主，避免把不同阶段的结构变化解释成某个知识点的升降趋势。
\clearpage
\subsection*{B.4 低年级决赛}
\noindent 填、答分别依原卷填空和解答题序。低年级卷中反复出现的 ODE（O）须单独补基础，不能仅按数学分析教程的章目准备。
{\small\renewcommand{\arraystretch}{1.18}
\begin{longtable}{P{2.8cm}P{1.8cm}P{10.8cm}}
\toprule 卷别 & 真题 PDF 页 & 各题主标签（依原题序）\\\midrule\endhead
5 低决 & 34--36 & 题序：G、L4、L3、O、A3、L5\\
6 低决 & 45--47 & 填：L4、A4、A3、L2\newline 答：G、L1、A2、O、A3\\
7 低决 & 56--58 & 填：L2、A4、A3、L1\newline 答：G、L4、A2、O、A3\\
8 低决 & 67--70 & 填：L5、A2、A3、L4\newline 答：G、L1、A4、A1、O\\
9 低决 & 79--81 & 填：L2、A1、A3、L4\newline 答：G、L1、A3、A4、O\\
10 低决 & 90--93 & 填：L4、A2、A5、L4\newline 答：G、L1、A2、A1、O\\
11 低决 & 106--109 & 填：A1、A3、G、L4\newline 答：G、L4、A3、A4、O\\
12 低决 & 121--125 & 填：A3、A1、L1、L1\newline 答：G、L3、A2、A1、O\\
13 低决 & 142--145 & 填：A1、A1、L1、O\newline 答：G、L3、A2、L4、A2、O\\
14 低决 & 167--170 & 填：A2、A3、L1、A3\newline 答：G、L3、A2、L1、A1、O\\
15 低决 & 183--186 & 填：L4、A1、A3、O\newline 答：G、L3、A3、L1、A1、O\\
16 低决 & 199--202 & 填：L4、A2、A3、O\newline 答：G、L5、A1、G、A1、O\\
17 低决 & 214--218 & 填：L1、A2、A3、O\newline 答：G、L1、A1、L4、A1、O\\
\bottomrule\end{longtable}}

\para{读表提示}
基础计算与结构证明同时存在。L3、L1、A1、A2 等标签之间常有交叉；例如矩阵函数既用代数关系，也用级数收敛。准备决赛时，先检查基本填空是否能稳定完成，再练跨主题证明。
\clearpage
\subsection*{B.5 高年级决赛}
\noindent 高低年级相同基础题保留原卷位置，不当作两次独立命题。R/C/H/D/P/N 为实分析、复分析、抽象代数、微分几何、概率和数值分析相关主标签。
{\small\renewcommand{\arraystretch}{1.18}
\begin{longtable}{P{2.8cm}P{1.8cm}P{10.8cm}}
\toprule 卷别 & 真题 PDF 页 & 各题主标签（依原题序）\\\midrule\endhead
5 高决 & 37--41 & 题序：G、L4、A3、L5、O、C、R、D、P、H、N\\
6 高决 & 48--52 & 填：L4、A4、A3、L2\newline 答：G、L1、A2、O、C、R、D、N、H、P\\
7 高决 & 59--63 & 填：L2、A4、A3、L1\newline 答：G、L4、A2、H、R、D、N、C、P\\
8 高决 & 71--75 & 填：L5、A2、A3、L4\newline 答：G、L1、A4、H、R、D、N、C、P\\
9 高决 & 82--86 & 填：L2、A1、A3、L4\newline 答：G、L1、A3、H、R、D、N、C、P\\
10 高决 & 94--99 & 填：L4、A2、A5、L4\newline 答：G、L1、A2、H、R、D、N、C、P\\
11 高决 & 110--115 & 填：A1、A3、G、L4\newline 答：G、L4、A3、R、C、H、D、P、N\\
12 高决 & 126--131 & 填：A3、A1、L1、L1\newline 答：G、L3、A2、R、C、H、D、P、N\\
13 高决 & 146--151 & 填：A1、A1、L1、O\newline 答：G、L3、A2、R、C、H、D、P、N\\
14 高决 & 171--176 & 填：A2、A3、L1、A3\newline 答：G、L3、A2、R、C、H、D、P、N\\
15 高决 & 187--192 & 填：L4、A1、A3、O\newline 答：G、L3、A3、R、C、H、D、P、N\\
16 高决 & 203--207 & 填：L4、A2、A3、O\newline 答：G、L5、A1、R、C、H、D、P、N\\
17 高决 & 219--225 & 填：L1、A2、A3、O\newline 答：G、L4、A1、R、C、H、D、P、N\\
\bottomrule\end{longtable}}

\para{读表提示}
样本的后期高年级卷有稳定的专业课模块。这里的“考得少”必须限定在初赛或某一教材内；如果目标是高年级决赛，仅精读这两册教程仍不够。N 也包括逼近论、求积等数值相关内容，并非每年都只考线性迭代算法。
